% ED26-BM-OHTR2012-17: qualified complete 2012 chapter, not a 2002 substitute.
\chapter{Judgmental Heuristics: The History of a Research Program}
\label{sec:p19-heuristics-program-history}

A probability judgment often arrives before the judge can explain how it
was made. The historical question is not simply whether it is correct, but
what operations produced it, why those operations usually provide useful
information, and why they sometimes produce a systematic departure from a
specified benchmark. Griffin, Gonzalez, Koehler and Gilovich organize their
2012 history around this connection between a normative comparison and a
positive psychological explanation
\cite[pp.~322--345]{GriffinEtAl2012HeuristicsHistory}.
Their account follows antecedents, three waves of the heuristics-and-biases
program, its major objections and its changing explanatory vocabulary.
It is a history written in 2012, not a claim that subsequent research has
settled every dispute it describes.

The empirical reports below retain that source boundary. Except for the
explicitly linked original conjunction and calibration records, earlier
studies are presented as this historical chapter reports them. Their
missing recruitment, exclusion, uncertainty and replication fields are
not reconstructed from remembered results. The distinction between a
historical account, an original experiment and an independently justified
mathematical benchmark follows \cref{def:p19-evidence-record}.

\section{The rational actor as benchmark, not a process diagram}
\label{sec:p19-hph-rational-actor}

The opening examples contrast a politician's rapidly produced odds about
Kosovo peace negotiations with a newspaper's extraordinary-odds treatment
of a family whose eight children alternated girl and boy. The first
illustrates the availability of an impression without an articulated
calculation; it supplies neither an outcome-frequency study nor evidence
that the odds were wrong. The second illustrates the confusion between a
specified sequence and a class of sequences judged to look random.
An orderly-looking birth pattern may be unusual as a visual pattern while
having the same probability as each other specified four-of-each order
under an appropriate independent-birth model. That model, not appearance
alone and not the probability axioms alone, supports the comparison.

The source's rational-actor framework combines subjective probabilities
and utilities with consistency constraints. It is not a contrast between
an emotional person and an unemotional person. Preferences determine
what an actor values; probabilities represent uncertainty; an expected-
utility comparison combines them. A coherent preference representation
is not itself a discovery of the actor's beliefs, and a probability
axiom is not a psychological algorithm for producing a judgment.
The bounded-rationality distinctions in
\cref{def:p17-br-principle-theory,def:p17-br-three-rationalities}
keep substantive criteria, feasible procedures and ecological success
separate.

\begin{hbdefinition}{def:p19-hph-two-evaluations}{Coherence and calibration in the history}
\emph{Coherence} compares related judgments under one interpretation and
information set. A probability assigned to an intersection cannot exceed
the probability assigned to an event containing it. \emph{Calibration}
compares judgments with outcomes under a specified judgment--outcome law:
among the positive-mass occasions receiving a probability of .70, the
conditional outcome probability should be .70. The finite definition is
\cref{def:p19-calibration-law}; event inclusion is treated in
\cref{thm:p19-conjunction-benchmark}.

A person's judgments may satisfy an inclusion comparison while being
miscalibrated, or may have a correct mean confidence while failing
conditional calibration. A finite table's observed frequencies do not
automatically equal the population law. Neither benchmark alone identifies
the procedure that generated a response.
\end{hbdefinition}

The historical account uses Bayes' rule to illustrate how a normative
updating model combines prior odds and the diagnostic value of evidence.
It also recognizes disagreement over the meaning of single-event
subjective probabilities versus repeatable frequencies. Such disagreement
does not license treating a similarity ranking as if it were a coherent
probability merely because both are expressed in numbers.

\begin{hbdefinition}{def:p19-hph-odds-domain}{A two-hypothesis updating comparison}
Let \(H_1,H_2\) partition a probability space, with
\(P(H_1)>0\), \(P(H_2)>0\). Let evidence \(D\) have positive
probability and let \(P(D\mid H_2)>0\). Write
\(o_0=P(H_1)/P(H_2)\) for prior odds and
\(\ell=P(D\mid H_1)/P(D\mid H_2)\) for the likelihood ratio.
Both likelihoods refer to the same evidence and the same hypothesized
sampling model. A likelihood ratio is not the posterior probability of
the hypothesis.
\end{hbdefinition}

\begin{hbtheorem}{thm:p19-hph-odds}{proved=annex:p19.hph.odds}{The source's prior-odds illustration}
Under \cref{def:p19-hph-odds-domain},
\[
 \frac{P(H_1\mid D)}{P(H_2\mid D)}=\ell o_0,
 \qquad P(H_1\mid D)=\frac{\ell o_0}{1+\ell o_0}.
\]
For prior odds \(1:9\), a likelihood ratio \(9:1\) gives posterior
odds \(1:1\), hence probability \(1/2\); a likelihood ratio \(1:9\)
gives odds \(1:81\), hence probability \(1/82\).
These are the source's stipulated updating examples, not measured
judgments or clinical risk estimates.
\end{hbtheorem}

The disease-versus-no-disease language of the historical illustration
serves only to interpret two hypotheses. No screening recommendation
follows without actual prevalences, test properties, costs and a decision
problem. In the history, the point is that identical evidence need not
produce identical posterior beliefs when priors differ, and that a strong
impression of evidence is not enough to determine its proper weight
\cite[pp.~323--324]{GriffinEtAl2012HeuristicsHistory}.

\section{Antecedents: limited computation and imperfect statistical judgment}
\label{sec:p19-hph-antecedents}

The chapter places the program in a postwar convergence. Engineering and
military problems made expert judgment important; the cognitive revolution
made information processing a psychological subject; computing supplied
both metaphors and candidate procedures; economics exposed the demands
of fully optimizing choice. These antecedents did not all offer the same
account of error.

Simon objected to models demanding exhaustive information and computation.
The history identifies three components in satisficing: search begins
with accessible or local alternatives, an aspiration criterion determines
when an acceptable option has been found, and future consequences receive
a coarse evaluation rather than an exhaustive optimization. Attention
and memory limits help explain why few alternatives are considered.
The positive proposal is that a feasible procedure can yield good
approximations in suitable environments; the negative claim is that
an unrestricted optimizer is not an adequate account of the actual
procedure. The source connects this approach with artificial-intelligence
heuristics, organizational search and later work on expertise. It does
not derive the heuristics-and-biases program merely by adding a resource
bound to Bayes' rule. The full aspiration/search distinction is developed
at \cref{def:p17-br-satisficing}
\cite[pp.~324--325]{GriffinEtAl2012HeuristicsHistory}.

Ward Edwards' updating experiments supplied a different starting point.
People observed evidence such as colored poker chips and revised odds.
The historical characterization is that their updates were often in the
Bayesian direction but too conservative: the new evidence received too
little weight. Important exceptions involved large shifts in response to
weakly probative evidence. A theory built only around underreaction
therefore had a restricted explanatory domain even before the new program.

Other antecedents included Adams and Adams on confidence in knowledge;
Pollack and Decker on confidence in difficult, noisy perceptual decisions;
and Cohen on psychological probability, chance and luck. The history
contrasts overconfidence on hard perceptual tasks with better calibration
or underconfidence on easier tasks. It also distinguishes choosing in a
random environment from reporting its probability. In probability
matching, a chooser distributes choices roughly like the outcome
frequencies even when always selecting the more common outcome would
yield more expected successes under the supplied conditions.

\begin{hblemma}{thm:p19-hph-patterns}{proved=annex:p19.hph.patterns}{Specified patterns, matching and ratio comparisons}
Suppose eight births are independent with a fixed girl probability
\(p\in(0,1)\). Each specified order containing four girls and four
boys has probability \(p^4(1-p)^4\). For \(p=1/2\), each is
\(1/256\); the class of all such orders has probability \(70/256\).
Conditioning on four of each gives every such order probability \(1/70\).
The event of one specified alternating order is not the class of all
balanced orders.

If independent trials produce red with probability \(q=.7\), and an
independent chooser selects red with probability \(a\), its one-trial
success probability is \(aq+(1-a)(1-q)\). Matching \(a=q\)
gives .58, whereas always selecting red gives .70. This comparison
uses expected successes as its objective.

For a uniformly sampled bean with the same reward for red and no other
cost difference, the red probabilities in urns with 1 red of 10,
9 red of 100, and 5 red of 100 are .10, .09, and .05. The first
maximizes expected reward among those three urns. Numbers of red beans
alone do not determine that comparison.
\end{hblemma}

The birth calculation does not assume actual births are independent and
identically distributed. Nor does the matching calculation prove that
every matching behavior is irrational: exploration, changing contingencies,
other objectives or dependencies change the problem. These conditions
make explicit what the historical toy examples require.

Clinical prediction was another antecedent, not a separate demonstration
that clinicians are always inferior to algorithms. The source reports
Meehl's comparisons of judgment and actuarial prediction, neglect of
base rates, Oskamp's finding that additional case information could
increase confidence without matching accuracy, and Mischel's criticism
of interview-based cross-situation prediction. These review reports
connect confidence, evidential validity and the attraction of individual
case information. They are not one pooled experiment or a current clinical
recommendation. The history provides no common sampling frame or effect
size by which to aggregate them
\cite[pp.~325--326]{GriffinEtAl2012HeuristicsHistory}.

\section{A methodological departure and three waves}
\label{sec:p19-hph-waves}

Tversky's mathematical psychology and Kahneman's work on perception
converged on a possibility different from a slightly noisy or conservative
Bayesian calculator: natural judgment might use qualitatively different
operations. The analogy was a perceptual illusion. The horizon moon and
the overhead moon need not look equally large because size perception
constructs an interpretation using distance cues. Likewise, using visual
clarity as a distance cue generally helps but misleads when visibility
conditions vary. An illusion isolates the operation by finding conditions
in which it gives a characteristic error.

This method targets a diagnostic failure, not a random sample of all
occasions on which perception or judgment succeeds. Statistical knowledge
can coexist with a compelling contrary appearance, as knowledge of a
visual illusion can coexist with its perceived effect. Gestalt ideas
about organizing partial information into a coherent whole help explain
why a complete numerical analysis need not precede an immediate judgment.
The program's negative message was that people need not be conservative
Bayesians; its positive message concerned the constructive operations
they use. Neither message was simply that people are unintelligent.
The history also connects later prospect theory with reference dependence,
diminishing sensitivity and asymmetric responses to gains and losses;
that distinct choice theory has its own presentation at
\cref{sec:p17-rationality-8-3}
\cite[pp.~326--328]{GriffinEtAl2012HeuristicsHistory}.

The first wave, located by the authors in 1969--1974, emphasized accessible
operations and striking demonstrations: similarity, ease of generating
instances, and starting points with adjustment. The second wave moved
toward the combination of information, causal interpretations, evidence
strength and weight, and models of comparative support. The third placed
a broader set of intuitive judgments within dual-process accounts.
The waves describe an evolution of explanations and methods; they are
not three disjoint populations or a finding that later labels replaced
all earlier mechanisms.

Early between-person demonstrations often used births, schools and
professions rather than explicitly random devices. This was deliberate:
chance devices can cue taught statistical routines and conceal the
natural assessment the experiment aims to investigate. The choice raises
an important design question about task interpretation and transfer.
The authors defend process-diagnostic demonstrations without claiming
that demonstration rates estimate everyday error prevalence. Their
contrast with Simon is therefore partly a contrast of investigative
purpose: finding conditions that expose an operation versus documenting
when a feasible procedure succeeds in its environment.

\section{First wave: three different operations}
\label{sec:p19-hph-first-wave}

\begin{hbdefinition}{def:p19-hph-three-heuristics}{Operations and the target of judgment}
\emph{Representativeness} assesses a case by its resemblance to a category,
outcome-generating process or causal pattern. \emph{Availability} assesses
frequency or possibility using the ease with which relevant instances
can be recalled or constructed. \emph{Anchoring and adjustment} begins
with an accessible value and moves away from it to estimate a target.

These descriptions name proposed operations. They do not imply that a
judge consciously selects a shortcut, that every task uses one operation,
or that a final biased response reveals its mechanism uniquely. The
comparison between a heuristic attribute and the requested target is
essential: similarity is not probability, retrieval ease is not frequency,
and an available starting value is not evidence about the target.
\end{hbdefinition}

\subsection{Representativeness: randomness, prediction and category resemblance}

An apparently random sequence is expected to look locally irregular.
This expectation can make a short sample seem more balanced than actual
random sampling warrants. The source connects this with the gambler's
fallacy, probability matching and its historical account of the basketball
hot-hand literature. Its basketball report describes the literature then
available; it is not a present-day theorem that shooting probabilities
never change or that later selection-sensitive analyses are irrelevant.
Randomness, stationarity and independence must be independently justified.
The fixed-string distinctions already developed in
\cref{thm:p19-r24-strings} prevent a pattern's representativeness from
substituting for its specified probability.

Prediction by evaluation offers a related route. An extreme profile can
produce an equally extreme forecast even when its predictive validity is
limited. Regression toward the mean then seems to demand a special causal
explanation, such as praise or punishment, when part of the movement is
already expected under an imperfect predictive relation. The history
connects parental extremes, replication expectations and the law of
small numbers: a small sample can be treated as if it faithfully displayed
the population's structure. Neither a sample's resemblance nor confidence
in it determines the sample size needed for reliable inference.

The Tom W demonstration makes the case/category comparison explicit.
A graduate-student profile emphasizes orderliness and competence together
with low creativity and weak social engagement. Different groups rank
nine fields by similarity to the profile or likelihood of enrollment;
another estimates how common those fields are. In the chapter's report,
aggregate likelihood and similarity rankings correspond closely, whereas
the base-rate ordering does not. These are different-group rank profiles,
not individually observed algorithms. The chapter does not supply the
full recruitment, exclusion or uncertainty record here.

Base-rate neglect is consequently not the claim that people can never use
a base rate. The source discusses order of presentation, a physical sampling
schema and causally meaningful rates as conditions encouraging its use.
The existence of an effect under one description is compatible with its
reduction under another. A causal base rate, such as failure on a difficult
exam, can shape the case's representation rather than compete with it as
an unrelated number. The historical anecdote that food shaped like a
heart, brain or kidney is good for that organ illustrates resemblance
being substituted for a causal mechanism; it supplies no nutritional or
medical evidence
\cite[pp.~329, 331--332]{GriffinEtAl2012HeuristicsHistory}.

\subsection{Availability: accessible content and experienced ease}

Ease of recall can be informative when the retrieval procedure and
experience are related to the target frequency. Thinking of Canadian
comedians may use genuine cultural exposure. But the available sample,
its salience, and the cue used to search memory need not track population
frequencies. A word beginning with \emph{r} is easier to retrieve than a
word with \emph{r} in the third position. The source reports that the latter
class was more frequent in the historical word task, about three times
as frequent. This is a report about that task and its lexicon, not a
claim about every language, corpus or current English dictionary.
The exposure-versus-retrieval distinction in
\cref{thm:p19-r24-exposure} is a separate reason not to equate what is
remembered with an unbiased census.

Ross and Sicoly's spouses estimated their own contributions to household
work, the social calendar and quarrels. The chapter reports sums exceeding
the total available contribution, including negatively valued activities.
That pattern supports an availability account over a wholly flattering
self-enhancement account: one's own acts are especially accessible even
when they are undesirable. It does not uniquely isolate availability;
differences in interpretation or access to the partner's behavior remain
possible. The source does not provide a participant-level reconstruction
from which a new interval or causal estimate can be derived.

Schwarz and colleagues' six-versus-twelve examples distinguish the amount
of recalled content from the felt ease of producing it. Respondents asked
for six examples of assertive behavior rated themselves more assertive
than those asked for twelve; the corresponding unassertive-example
comparison went the other way. In the chapter's account, difficulty
producing twelve assertive examples made assertion seem less characteristic,
while difficulty producing twelve unassertive examples made it seem more
characteristic. The manipulation does not merely show that more recalled
instances always raise a judgment. The chapter also reports Fox's classroom
example: requesting ten improvements rather than two could yield a better
course evaluation, consistent with difficulty generating many criticisms.
No measured effect size, exclusion rule or confidence interval is supplied
by this historical summary, and course ratings are not a direct measure
of teaching quality
\cite[pp.~330--331]{GriffinEtAl2012HeuristicsHistory}.

\subsection{Anchoring: a family rather than one universal adjustment process}

For a self-generated anchor, a person asked about Venus's year may begin
with Earth's 365 days and adjust downward because Venus is nearer the Sun.
The source's illustrative 275-day answer remains above its stated 224.7-day
target. This example makes the direction of adjustment intelligible; it
is not a sampled estimate of how often people make that answer. The classic
externally supplied anchor is different: a wheel-generated comparison
number shifts estimates of the percentage of African members in the UN.
The chapter's question is why a value with no informational entitlement
can influence the later answer.

The history separates at least three accounts. For self-generated anchors,
insufficient adjustment can be a consciously reported process. Experimental
conditions involving distraction or a compatible head movement, and weak
accuracy motivation, can change adjustment. Experimenter-provided anchors
need not show the same pattern. Selective accessibility instead proposes
that comparing the target with an anchor activates anchor-consistent
knowledge: a very long versus short Nile comparison, or an expensive-car
comparison followed by recognition of luxury-car names, directs accessible
content toward the comparison. A third family involves numerical priming
without the same target-specific semantic relation.

The source includes cross-modal demonstrations: a long or short imagined
Mississippi represented by toothpicks or completion of long words influences
a later numerical estimate. It also contrasts numerically large 7300
meters with equivalent 7.3 kilometers in an unrelated project-cost task,
and telephone-model labels P97/P17 with later European sales estimates.
These comparisons motivate a distinction between semantic magnitude and
the priming of a number. They do not establish that all anchoring has
one mechanism, nor that semantic priming alone explains every result.
The chapter cites a 2010 exception to a broad claim that incentives never
affect adjustment from provided anchors. Mechanism-specific task differences
matter more than treating every anchor effect as interchangeable
\cite[pp.~332--334]{GriffinEtAl2012HeuristicsHistory}.

\section{Objections that changed the program's questions}
\label{sec:p19-hph-objections}

\subsection{Selective demonstrations, incentives and transfer}

An existence demonstration establishes that an effect can occur under
its instructions, information and response measure. It does not estimate
the proportion of all everyday judgments made through that operation.
The history explicitly resists a rhetorical move from an illusion-like
demonstration to a global verdict on human competence. Conversely, finding
conditions in which a bias diminishes does not refute its occurrence in
the original conditions. A useful objection identifies which assumption,
interpretation or process a comparison actually tests.

Information-integration approaches, including Anderson's cognitive
algebra, investigated how multiple sources are combined. Showing that
people can use base-rate information in a within-person integration task
constrains the claim that they lack the capacity altogether. It does not
show that the same weighting is used spontaneously in a different
between-person task. Capacity, application, task relevance and the
interpretation of the supplied information remain distinct issues.

Economic objections ask whether incentives, feedback, repetition or
markets remove effects found in brief laboratory judgments. In the
history's report, Grether's transparent bingo tasks yielded some improvement
with incentives while preserving departures; Camerer's repeated market
tasks showed attenuation with experience. These are different designs,
not one aggregate treatment effect of money. The chapter also notes
that discovering a base rate need not guarantee using it in a later case
judgment. These examples motivate rather than settle questions about
training, feedback and organizational correction. The selected review
does not supply standardized effect sizes, participant exclusion records
or sampling uncertainty for a new quantitative synthesis
\cite[pp.~334--335]{GriffinEtAl2012HeuristicsHistory}.

\subsection{Conversation and the meaning of the requested judgment}

The conversational objection has a constructive and a stronger form.
The constructive form says that people interpret relevance, wording and
the investigator's purpose in context. This agrees with a psychology of
active interpretation and can explain when a statistic is treated as
useful. The stronger form treats the apparent bias as an artifact and
rejects its transfer to judgments outside that conversation. These
positions cannot be collapsed into the statement that wording matters.
One must identify what interpretation the alternative procedure encourages
and whether the comparison still tests the intended probability relation.

Macchi's example in the history concerns an ambiguous conditional statistic
about married or single people and suicide. Rewording to clarify which
population is conditioned upon increased base-rate use. The chapter's
historical numbers are task materials, not current mortality rates.
The lesson is directional: the probability of a demographic attribute
among a selected outcome group is not the probability of that outcome
among the demographic group. A helpful rewording can remove this ambiguity
without demonstrating that all earlier effects arose from it. Endless
variants of a question are uninformative unless their change targets
a specified interpretation or mechanism.

The source also considers a scientific versus clinical perspective and
the description of information as randomly computer-selected versus
psychologist-selected. Such manipulations change the relevance a judge
may attribute to individuating information. The authors warn that a
single low-engineer-base-rate condition cannot establish proper base-rate
use: without a contrasting base-rate condition, an estimate's proximity
to one benchmark does not reveal how that statistic affected it. The
comparison concerns a changed construal of evidence, not simply an
experimenter making the same message clearer
\cite[pp.~334--335]{GriffinEtAl2012HeuristicsHistory}.

\subsection{Frequencies, single events and unequal comparison tasks}

A separate challenge contrasts subjective probabilities of unique events
with frequencies in a class. The history answers that subjective probability
still has coherence constraints, while recognizing that frequencies can
make those constraints easier to see. Its reference to a monetary-loss
argument for incoherent odds explains a normative motivation; it is not
a complete Dutch-book theorem supplied by this chapter, and no new
formal result is claimed from that quotation.

The authors emphasize that the heuristics-and-biases literature itself
recognized frequency advantages. A concrete set of cases can make a
subset relation transparent, supply background information and encourage
unpacking into component events. A holistic case description may instead
call up an immediate impression. Those psychological differences can be
real without proving that only frequency judgments are meaningful.

The dispute therefore concerns why formats help and what has been held
constant. Some critics propose an evolved specialization for frequencies;
the historical authors emphasize the accessibility of extensional reasoning,
concrete set inclusion and presentation. They point to difficulties in
claims that subjective chance is an unfamiliar invention and to the use
of probability language by statistically trained respondents. These are
arguments within the 2012 controversy, not independent evidence here about
the evolution of a cognitive module. Nonparallel response scales,
within-person versus between-person comparisons, and changes in task
content confound a simple probability-versus-frequency contrast. Making
one condition transparent while leaving the other opaque does not isolate
the format label
\cite[pp.~335--336]{GriffinEtAl2012HeuristicsHistory}.

\section{Second wave: multiple sources and the structure of evidence}
\label{sec:p19-hph-second-wave}

\subsection{Causal representation and inside/outside judgments}

In the second wave the question shifted from identifying a shortcut to
explaining which information enters an impression and how it is combined.
Causal base rates can be natural inputs to a representation. A difficult
examination explains a high failure rate; repeated success by a team can
be read as a propensity rather than an unrelated statistic. The chapter
also describes second-order representativeness: a sample may be judged
by whether its composition resembles its population. Statistical
information is consequently not psychologically uniform. The way it
is connected to the case can change its use.

The inside/outside planning contrast is another example. An inside view
constructs the project's particular sequence, difficulties and intentions;
an outside view considers the outcomes of an appropriate reference class.
The source locates the planning-fallacy account in this period. It does
not claim that the outside view ignores all case information or that any
available reference class is relevant. The important contrast is between
a holistic case-based construction and an analytic comparison, not an
invariant opposition between descriptive and numerical information.

The chapter distinguishes immediate natural assessments from broader
question-answering operations involving similarity, causal potency,
counterfactual comparison and surprise. A counterfactual account can
make a result seem unusually probable or surprising without measuring
its objective frequency. This broader view also allows biased judgment
to result from too much case-based construction, not merely from an
unmotivated cognitive miser using too little thought. Mental contamination
names a failure to keep an initially generated interpretation from
influencing a later judgment; it is not identified by an error score alone
\cite[pp.~336--338]{GriffinEtAl2012HeuristicsHistory}.

\subsection{The conjunction paper's exact place in the corpus}
\label{sec:p19-hph-conjunction-place}

The 1983 conjunction work is central to the second wave in this history,
not merely an early illustration of a memorably named bias. It tests
when adding a representative, easily generated or causally connected
condition makes a narrower description feel more probable. The chapter
uses Bill, a profile favoring an accountant stereotype over a jazz hobby,
as an example: adding accountant to jazz can improve descriptive fit
while reducing the event's extension.

The original-article treatment at
\cref{sec:p19-conjunction-persons,tab:p19-conjunction-ranks}
keeps the comparison's operational details. In its naive direct-ranking
Bill row, 94 Stanford/UBC undergraduate respondents without
probability/statistics training supplied likelihood ranks; the article
reports 92\% placing accountant-and-jazz above jazz, with common-filler
mean ranks 2.5 and 4.5. Those numbers are ranks and a reversal proportion,
not cardinal probabilities or a general-population error estimate.
The linked original record identifies its indirect stage, the educational
groups and the fact that the two formats include repeated contributions
by participants. The historical chapter's isolated 92\% is therefore
not a second independent sample
\cite[p.~298, table~1]{TverskyKahneman1983Conjunction}
\cite[p.~337]{GriffinEtAl2012HeuristicsHistory}.

The comparison with \cref{thm:p19-conjunction-benchmark} requires a common
event space, a common information set and the intended inclusive meaning
of jazz. A respondent answering whether the profile is typical, or
interpreting the unqualified hobby as excluding the occupation, answers
a different question. Such possibilities require evidence about
interpretation. They are not settled solely by observing a reversal.

The full original corpus is wider than Bill and Linda. Availability-based
word searches, clinical and sporting judgments, causal conjunctions,
betting and frequency formulations occupy their actual sections in
\cref{sec:p19-conjunction-judgments}. The source history stresses
nontransparent versus transparent comparisons, language controls and
frequency judgments, including the controversy and later adversarial
collaboration over format effects. It acknowledges that making inclusion
visible can reduce reversals while maintaining that many nontransparent
judgments follow the compelling case representation. Knowledge of a
rule and a contrary initial feeling can coexist. This is evidence about
the conditions of rule use, not proof that the participant lacks the rule.
The already-written task families are neither silently duplicated nor
credited here as a newly acquired original corpus.

\subsection{Strength and weight, confidence and calibration}
\label{sec:p19-hph-confidence-place}

Strength describes the apparent force or direction of the evidence;
weight concerns factors such as validity and sample size that determine
how much the impression warrants. The strength--weight account has a
judge begin with the impression's strength and adjust for its weight.
Underweighting validity or sample size can yield overconfidence when
weight is low and underconfidence when it is high. Thus the account
aims to explain both departures, unlike a single conservatism parameter.
These are explanatory predictions conditional on the evidence situation,
not a theorem that every small sample produces overconfidence.

The historical chapter mentions confidence findings in physicians,
economists, lawyers and weather forecasters. These are multiple reviewed
settings, not a common measured population. The relevant corpus location
is \cref{sec:p19-calibration-extreme-confidence}, where the original
Fischhoff--Slovic--Lichtenstein experiments retain task formats, recruited
volunteers, odds thresholds, repeated judgments, answer keys and printed
inconsistencies. Their original report's confidence proposals remain
explanations needing further process evidence, not a uniquely established
mechanism. Brier's scoring benchmark and the handbook's complete
calibration proofs occupy that chapter's distinct scoring sections.

In particular, answer accuracy is not calibration: a harder item set
can change correctness and the confidence distribution simultaneously.
Mean confidence, the conditional outcome rate at one confidence value,
resolution across confidence classes and expected squared error are
different quantities. Changing their sampling mixture can change a
calibration curve without demonstrating a change in one cognitive
operation. The history's strength--weight account provides a process
proposal for these patterns; the joint-law identities in
\cref{def:p19-calibration-law,thm:p19-calibration-decomposition}
do not prove that proposal
\cite[pp.~334, 338]{GriffinEtAl2012HeuristicsHistory}.

\subsection{Support theory: a description-sensitive probability model}

Support theory makes descriptions, rather than only extensional events,
the objects of comparative judgment. A hypothesis's support can come
from similarity, ease of generating instances or base-rate information.
Different descriptions of the same event can elicit different support.
This is a descriptive hypothesis about judgments; probability theory
does not require the probability of an event to change when it is renamed.

\begin{hbdefinition}{def:p19-hph-support}{Support for compared descriptions}
For compared descriptions \(A,B\), let \(s(A),s(B)\) be nonnegative
support scores with \(s(A)+s(B)>0\). The source model's judged
probability of \(A\) relative to \(B\), retaining the source's
two-description syntax \(P(A,B)\), is
\[
             P(A,B)=\frac{s(A)}{s(A)+s(B)}.
\]
Here \(A,B\) denote descriptions, not necessarily an arbitrary pair of
events in a sigma algebra. If the task interprets them as complementary
hypotheses, the relative comparison is a judgment about that task.
Here the two-argument \(P\) has a locally declared judgment-function
type, unlike the event probability law in the earlier odds comparison.
The scores are descriptive model inputs, not observed population
probabilities. The ratio alone imposes no relation between the support
of a packed description and support under its unpacking.
\end{hbdefinition}

\begin{hblemma}{thm:p19-hph-support}{proved=annex:p19.hph.support}{Normalization and the direction of unpacking}
For nonnegative scores with positive sum, \(P(A,B)\in[0,1]\) and
\(P(A,B)+P(B,A)=1\). Keeping \(s(B)>0\) fixed, an increase
in \(s(A)\) strictly increases \(P(A,B)\); keeping \(s(A)>0\)
fixed, an increase in \(s(B)\) strictly decreases it. If the other
score is zero, the corresponding ratio can remain at an endpoint.
With packed scores \(s(A)=s(B)=2\), raising only \(s(A)\) to 4
gives \(2/3\) instead of \(1/2\). This is constructed arithmetic,
not fitted experimental data.
\end{hblemma}

The source's Jack example compares lawyer with engineer. Unpacking lawyer
into corporate, criminal, divorce and tax specializations may activate
additional supporting possibilities and raise its judged relative
probability; unpacking engineer instead may lower the lawyer judgment.
The direction needs a psychological premise that unpacking increases
support. Implausible component descriptions can do the opposite, so
subadditivity is not a consequence of the normalized-ratio formula alone.
The source discusses unpacking/repacking and random-support developments
as attempts to specify these relationships.

Good agreement between judgments and proxy support measures makes a
model empirically useful under the tested tasks. It does not establish
that the proxy is the unique mental mechanism, that all fitted parameters
transfer to new tasks, or that the resulting judgments satisfy extensional
additivity across all descriptions. This formalization advances beyond
merely naming a heuristic while leaving specific psychological and
predictive obligations visible
\cite[p.~338]{GriffinEtAl2012HeuristicsHistory}.

\section{Broader second-wave assessments: affect and fluency}
\label{sec:p19-hph-affect-fluency}

\subsection{Affect and fluency broaden the available attributes}

An affective reaction can supply a readily available good/bad evaluation
when a person judges risk, value, agreement or whether to approach an
object. The affect heuristic therefore broadens the proposed inputs to
judgment beyond remembered counts and similarity. It is not a claim
that every use of emotion is irrational: a reaction may encode relevant
experience, and its adequacy depends on the target and environment.

Fluency concerns processing ease, which differs from the ease of
generating examples in an availability task. Imagining a future election
or illness can make its scenario easier to process and hence seem more
likely. Familiarity, category membership and the apparent validity of a
rhyming adage can likewise be influenced by processing ease. The history
also discusses disfluency as a possible cue to slow down and proceed
cautiously; in the reported difficult product-choice task, disfluency
increased selection of a default option. It therefore does not imply
a universal departure from defaults. Perceived psychological distance
is another application. A difficult-to-read city name can feel farther away without
changing geographical distance. These are the chapter's reviewed
applications; they are not a universal prescription to improve decisions
by making text illegible
\cite[p.~339]{GriffinEtAl2012HeuristicsHistory}.

\section{Third wave: automatic impressions and reflective intervention}
\label{sec:p19-hph-third-wave}

\subsection{Dual-process accounts and attribute substitution}

The two-process vocabulary contrasts relatively rapid, associative,
automatic processing with more deliberate, rule-based processing. The
history traces antecedents to Neisser and Wason/Evans, rather than treating
the vocabulary as a discovery that appeared fully formed in a popular
book. Kahneman's attention work already addressed information processed
without voluntary selection, and the 1983 argument implicitly permitted
a rule-governed route alongside a natural assessment. Later cognitive
and social models made that contrast explicit.

The source's jelly-bean example dramatizes a conflict between an understood
ratio and an attractive count. In the reviewed task, choosing the larger
bowl offered 9 red beans among 100 rather than 1 among 10; the chapter
reports 61\% choosing it. When the large bowl had only 5 red among 100,
23\% still chose it. The measures are choices between the supplied bowls,
not errors in verbally calculating their proportions and not a moral or
universal intelligence score. The mathematical comparison is the
equal-reward, uniform-draw part of \cref{thm:p19-hph-patterns}.
The historical summary does not supply participant count, full recruitment,
exclusions, uncertainty intervals or all procedural details, so these
percentages cannot become a new original-study population estimate here.
Knowing that .10 exceeds .09 or .05 is compatible with feeling attracted
to the greater number of favorable beans
\cite[pp.~339--340]{GriffinEtAl2012HeuristicsHistory}.

Attribute substitution asks how an accessible answer comes to stand in
for the requested one. A difficult target attribute, such as likelihood
of a cause of death, can be replaced by an easier attribute, such as
the accessibility of examples. A subject can then answer promptly without
explicitly recognizing the replacement. The theory identifies a potential
source of the initial response; explaining its persistence requires a
second question about why checking fails to detect or correct it.
Calling the initial answer System 1 and the correction System 2 is not
yet a complete account of either process.

Training may sometimes enable detection of a misleading cue, and deliberate
analysis can revise a first answer. It can also reinforce a mistaken case
construction or fail to begin. Thus the initial heuristic need not be
the sole or final basis of judgment. The history extends this perspective
to happiness, retrospective pain with duration neglect, and willingness
to pay for environmental goods. These differ in their target constructs,
time scales and response measures; one cannot pool their reported
responses into a single measured propensity to be biased. A predictive
mechanism must specify which attribute is generated, when it is available,
what intervention can alter it, and how the final task response is scored
\cite[pp.~340--341]{GriffinEtAl2012HeuristicsHistory}.

\subsection{Competing architectures and a limited explanatory vocabulary}

Two-process accounts do not agree about the sequence. A default-
intervention architecture has one response arise first and another process
sometimes correct it. A parallel-competition architecture permits more
than one response to develop together. A verbal contrast between fast
and slow does not distinguish those architectures. The history includes
skeptical arguments that the two-system division groups heterogeneous
processes, and Stanovich's distinction between ability and thinking
dispositions within the reflective side. Its tripartite vocabulary
does not by itself supply a uniquely tested algorithm either.

The questions are empirical: which operations the manipulation changes,
when responses diverge, and what independent evidence distinguishes
correction from competition or from changed task interpretation. Response
speed, a final accuracy score or an association with a cognitive test
is insufficient in isolation. The caution here agrees with the score/
construct and causal-comparison distinctions of
\cref{def:p19-empirical-question,def:p19-evidence-record}.
The 2012 history calls dual-process ideas an organizing development,
not the completion of a fully specified general theory
\cite[p.~341]{GriffinEtAl2012HeuristicsHistory}.

\section{What the history leaves open}
\label{sec:p19-hph-open-questions}

The program's evolution has several continuities. It begins with a
benchmark but seeks a psychological operation; a diagnostic failure is
informative without representing every normal judgment; and a useful
operation can produce bias when its cue and target come apart. Later
models add combination rules, causal and counterfactual representation,
affect, fluency and intervention rather than proving that the original
three operations exhaust cognition.

The source closes by identifying expanding investigations of semantic
and counterfactual construction, neural evidence, experience versus
description and expert judgment. It formulates eight outstanding issues.
First, how do biases relate to general cognitive ability? Second, what
additional individual differences matter? Third, when are a biased
response and its need for correction detected? Fourth, does overriding
an intuition use the same inhibitory resources as Stroop, flanker or
go/no-go performance? Fifth, under what experience conditions do biases
diminish? Sixth, do organizations and markets attenuate them, and by
which procedures? Seventh, can anchoring unify other judgments as a
master heuristic, or does its mechanism diversity defeat that proposal?
Eighth, what are the positive and negative contributions of this program
to a general theory of psychology? These are stated as questions, not
as established answers
\cite[p.~341]{GriffinEtAl2012HeuristicsHistory}.

\begin{hbmethod}{def:p19-hph-historical-inference}{Using the history without erasing its evidence boundaries}
To evaluate one of the historical comparisons, identify its requested
target, the allegedly substituted attribute and the supplied benchmark.
Keep the observation's unit and measure fixed: a rank, a self-rating,
an urn choice and a probability report are not interchangeable.
Locate whether the report is an original article record or the 2012
authors' review of that record. Retain unavailable recruitment, exclusions,
uncertainty and replication information as unavailable rather than
filling a standard-looking table with invented values.

Then compare explanations at their actual point of disagreement.
Retrieval content versus experienced ease requires different manipulation
predictions; self-generated adjustment versus selective accessibility
requires different anchor conditions; causal relevance versus simple
base-rate use requires a genuine base-rate contrast; frequency format
versus inclusion transparency requires parallel tasks. Transfer to a
new population, environment or task is a separate claim. Neither
mathematical coherence nor an empirical departure alone establishes
that transfer or uniquely identifies a cognitive mechanism.
\end{hbmethod}

The distinction between psychological explanation and normative assessment
therefore survives the program's changes. Expected utility, Bayesian
updating, inclusion and calibration can specify comparisons under their
assumptions without describing how a person thinks. A measured operation
can explain an answer without deciding that the answer is best for an
unrelated objective. Criticism can refine the task or process account
without abolishing a demonstrated departure, and a successful task can
show a capacity without proving its universal deployment. This is the
source-specific lesson of the three-wave history, not a generic verdict
for or against the use of heuristics.
